% !TeX root = main.tex
代码能在自己的机器上按 F5 跑起来，和它能被装到另一台机器上跑起来，是两件完全不同的事。发布阶段的事故几乎都不发生在源码里，而是藏在配置切换、运行库选择、依赖收集和安装包这几个环节。这一章按“从 Release 配置到交付安装包”的顺序把这些环节逐个拆开，重点放在每一步为什么会导致崩溃、以及怎么提前验证。
\section{发布前的配置切换}
\subsection{Debug 与 Release 的本质差异}
切换配置不是换个名字，而是同时改掉了一整组编译器和链接器选项。默认两个配置的差异大致如下：

\begin{longtable}{L{2.8cm}L{5.8cm}L{5.9cm}}
\toprule
\textbf{维度} & \textbf{Debug 配置} & \textbf{Release 配置} \\
\midrule
\endhead
\bottomrule
\endlastfoot
优化级别 & \code{/Od} 关闭优化，并叠加 \code{/RTC1} 基本运行时检查 & \code{/O2} 最大化速度 \\
预处理宏 & 定义 \code{_DEBUG}，不定义 \code{NDEBUG} & 定义 \code{NDEBUG}，不定义 \code{_DEBUG} \\
\code{assert} & 生效 & 被 \code{NDEBUG} 整体去掉 \\
运行库 & \code{/MDd} 或 \code{/MTd} & \code{/MD} 或 \code{/MT} \\
迭代器调试级别 & \code{_ITERATOR_DEBUG_LEVEL=2} & \code{_ITERATOR_DEBUG_LEVEL=0} \\
调试信息 & \code{/Zi}，链接器默认生成 PDB & 默认不生成 PDB \\
增量链接 & 打开（\code{/INCREMENTAL}） & 关闭，并启用 \code{/OPT:REF}、\code{/OPT:ICF} \\
全程序优化 & 关闭 & 项目模板常在 Release 下打开（\code{/GL} + \code{/LTCG}） \\
未初始化内存 & 调试堆会填充特征字节，容易暴露 & 内容随机，可能出现“在我机器上是好的” \\
速度与体积 & 慢、大 & 快、小 \\
\end{longtable}

\begin{warnbox}[注意]
\code{assert} 在 Release 下是\textbf{完全不执行}的，所以不要把带副作用的表达式写进 \code{assert}（例如 \code{assert(p = CreateBuffer())}）。Debug 下它正常工作，Release 下 \code{p} 会永远是空指针，而且不会有任何提示。
\end{warnbox}

\begin{tipbox}[提示]
想保留 Release 的性能又想能调试，可以克隆一个中间配置：\code{/O2} + \code{/Zi}，并在 \ui{项目 > 属性 > 链接器 > 调试 > 生成调试信息} 选择 \textbf{是 (/DEBUG)}。如果问题只在优化后出现，就单独建一个把 \ui{项目 > 属性 > C/C++ > 优化 > 优化} 改成 \textbf{已禁用 (/Od)} 的配置，用二分法缩小范围。
\end{tipbox}
\subsection{常用优化开关的含义与风险}
\begin{longtable}{L{2.6cm}L{5.4cm}L{6.6cm}}
\toprule
\textbf{开关} & \textbf{含义} & \textbf{需要注意的风险} \\
\midrule
\endhead
\bottomrule
\endlastfoot
\code{/O2} & 最大化速度，等于打开一揽子子优化 & 行号会错位、局部变量可能被优化掉，单步调试体验骤降 \\
\code{/Od} & 关闭所有优化 & 性能损失明显，只用于排查阶段 \\
\code{/O1} & 最小体积，与 \code{/O2} 互斥 & 体积优先时速度可能反而下降 \\
\code{/Ob2} & 内联展开（\code{/Ob0} 禁止、\code{/Ob1} 仅 \code{inline}、\code{/Ob2} 全部） & 断点可能进不去被内联的函数 \\
\code{/Oi} & 生成内建函数（intrinsic） & 部分 CRT 调用不再可断点 \\
\code{/Oy} & 省略帧指针 & 调用栈回溯可能不完整，是“看不到局部变量”的主因 \\
\code{/GF} & 字符串字面量池化 & 相同字面量共享同一地址，用地址比较字符串会出错 \\
\code{/Gy} & 函数级链接 & 配合 \code{/OPT:REF} 剔除未引用函数，是减体积的关键 \\
\code{/GL} & 全程序优化，编译期只产出中间表示 & 链接明显变慢、内存占用高，必须由 \code{/LTCG} 完成最终代码生成 \\
\code{/LTCG} & 链接时代码生成 & 与增量链接冲突，链接器会提示忽略 \code{/INCREMENTAL}（\code{LNK4075}） \\
\code{/arch:AVX2} & 允许编译器生成 AVX2（256 位）指令 & 在不支持 AVX2 的 CPU 上会触发非法指令异常（\code{0xC000001D}） \\
\code{/arch:AVX512} & 允许生成 AVX-512 指令 & 支持该指令集的消费级 CPU 更少，风险比 \code{/arch:AVX2} 更大 \\
\code{/fp:fast} & 浮点快速模式，允许重排与近似 & 结果与调试版不一致，可能出现 \code{a + b != b + a}，故障无法复现 \\
\code{/GS} & 缓冲区安全检查，默认开启 & 关掉会放大栈溢出风险，不要关 \\
\code{/guard:cf} & 控制流保护 & 与某些加壳、注入类工具冲突 \\
\code{/W4}、\code{/permissive-} & 打开第 4 级警告、关闭宽松模式以贴近标准 & 会让老代码冒出一批新警告，建议在发布前逐个处理 \\
\end{longtable}

\textbf{关于 \code{/arch} 的正确用法}：混用不同 \code{/arch} 编译的静态库不会链接失败，但最终程序会要求运行机器支持其中最高的那套指令集。如果必须兼顾老机器，做法是用 \code{/arch:SSE2} 或 \code{/arch:AVX2} 出通用版本，再在运行时用 \code{__cpuid} / \code{__cpuidex} 读取 CPUID 特性位做派发。

\begin{keybox}[实战场景]
开发机是最新一代 CPU，用 \code{/arch:AVX2} 编译后自测一切正常。程序发到一台老笔记本上双击就退出，事件查看器里是异常代码 \code{0xC000001D}（STATUS\_ILLEGAL\_INSTRUCTION）。这就是把编译期指令集当成运行期保证的后果。
\end{keybox}
\subsection{切换到 Release 后先做的三件事}
\begin{enumerate}
  \item 打开 \ui{生成 > 配置管理器}，确认“活动解决方案配置”是 Release，并且列表里\textbf{每一个项目}都是 Release——依赖库项目漏切一个，整个发布就是 Debug 版 CRT 混进 Release 二进制。
  \item 打开 \ui{项目 > 属性 > C/C++ > 命令行}（右上角可以切换配置与平台），肉眼扫一遍真实命令行，确认没有从 Debug 继承下来的 \code{/Od}、\code{/RTC1}、\code{/MDd}。
  \item 执行 \ui{生成 > 重新生成解决方案}，而不是“生成”。增量生成会把 Debug 的 \code{.obj} 留在中间目录里，改配置后必须全量重来。
\end{enumerate}
\section{运行库选择：\code{/MD} 还是 \code{/MT}}
\subsection{四个取值}
位置在 \ui{项目 > 属性 > C/C++ > 代码生成 > 运行库}：

\begin{longtable}{L{3.6cm}L{2.4cm}L{8.6cm}}
\toprule
\textbf{选项} & \textbf{开关} & \textbf{说明} \\
\midrule
多线程调试 DLL & \code{/MDd} & Debug 版动态 CRT，带调试堆 \\
多线程调试 & \code{/MTd} & Debug 版静态 CRT \\
多线程 DLL & \code{/MD} & Release 版动态 CRT，依赖可再发行程序包 \\
多线程 & \code{/MT} & Release 版静态 CRT，不依赖外部 DLL \\
\bottomrule
\end{longtable}

选 \code{/MD} 时，程序在运行期需要下列文件：

\begin{longtable}{L{4.6cm}L{4.4cm}L{5.6cm}}
\toprule
\textbf{文件} & \textbf{作用} & \textbf{备注} \\
\midrule
\code{MSVCP140.dll} & C++ 标准库实现 & 由可再发行程序包安装 \\
\code{VCRUNTIME140.dll} & 编译器支持例程（异常、RTTI 等） & 同上 \\
\code{VCRUNTIME140_1.dll} & 较新工具集补充的异常处理支持 & 只装过旧的运行库就会缺这个文件 \\
\code{ucrtbase.dll}、\code{api-ms-win-crt-*.dll} & 通用 C 运行时（UCRT） & Windows 10 起属于系统组件，一般无需分发 \\
\code{concrt140.dll} & 并发运行时 & 只在用到并发容器等设施时才会加载 \\
\bottomrule
\end{longtable}
\subsection{取舍表}
\begin{longtable}{L{2.6cm}L{6.0cm}L{6.0cm}}
\toprule
\textbf{维度} & \textbf{\code{/MD}（动态 CRT）} & \textbf{\code{/MT}（静态 CRT）} \\
\midrule
产物体积 & exe 小，逻辑集中在系统与运行库目录 & exe 明显变大，通常增加数百 KB \\
部署成本 & 必须确保目标机装有对应位数的可再发行程序包 & 免依赖，拷贝即用 \\
跨模块传 C++ 对象 & 全进程共用一份 CRT，跨 DLL 传递 \code{std::string} 等相对安全 & 每个模块各有一份 CRT 状态，跨边界释放内存极易堆损坏 \\
安全更新 & 随运行库更新一起被修补 & 必须自己重新编译并重新发布 \\
第三方库匹配 & 绝大多数官方预编译包默认是 \code{/MD} & 需要第三方库也提供 \code{/MT} 版本 \\
典型场景 & 常规桌面程序、插件式架构、多模块工程 & 绿色单文件工具、无管理员权限的受限环境 \\
\bottomrule
\end{longtable}
\subsection{混用 \code{/MT} 与 \code{/MD} 的后果}
同一个进程里同时出现静态 CRT 和动态 CRT，等于有两套完全独立的运行时状态：两个堆、两份 \code{errno}、两份 \code{std::string} 的实现、两份异常处理机制。在 A 模块里 \code{new}、在 B 模块里 \code{delete}（或 \code{malloc} / \code{free} 交叉），就会踩坏堆。
\begin{itemize}
  \item 静态库混用时，链接器通常能用 \code{LNK2038} 拦住你，提示“检测到 RuntimeLibrary 的不匹配项：值 MT\_StaticRelease 不匹配值 MD\_DynamicRelease”。
  \item \code{_ITERATOR_DEBUG_LEVEL} 不一致（Debug 2 与 Release 0）也会报 \code{LNK2038}。
  \item \textbf{但如果配置不一致的是一份 DLL，链接期根本发现不了}，只能在运行期炸成堆损坏（\code{0xC0000374}）或随机崩溃。
\end{itemize}

\begin{warnbox}[注意]
第三方库的 CRT 必须与本项目一致。拿到一份 \code{.lib} 或 \code{.dll} 时，第一件事是确认它是 \code{/MD} 还是 \code{/MT}、是 Debug 还是 Release、是 x86 还是 x64。四个维度都对上，才能避免两套运行时的灾难。
\end{warnbox}

统一的工程约定：
\begin{enumerate}
  \item 一个解决方案内的所有项目统一使用同一种运行库，写进项目模板或 \code{.props} 属性表。
  \item 先看第三方库给的是什么：Qt、OpenCV、Boost 的官方二进制包基本都是 \code{/MD}，那就全工程 \code{/MD}。
  \item 如果第三方只提供 \code{/MT} 版本，要么全工程改 \code{/MT}，要么自己用 \code{/MD} 重新编译那个库。
  \item 用 \code{/MT} 时不要跨 DLL 边界传递 STL 容器与字符串，改用 C 风格接口或纯虚接口。
  \item 需要同时支持多模块时，优先 \code{/MD}，把“体积小、可被系统更新”的优势留给运行库。
\end{enumerate}
\section{静态链接第三方库与动态链接 DLL}
\subsection{两种方式的部署代价}
\begin{longtable}{L{2.6cm}L{6.0cm}L{6.0cm}}
\toprule
\textbf{对比项} & \textbf{静态库（\code{.lib}，静态链接）} & \textbf{动态库（\code{.lib} 导入库 + \code{.dll}）} \\
\midrule
需要拿到 & \code{.lib} 与头文件 & 导入库 \code{.lib}、头文件、以及 \code{.dll} \\
部署时要带什么 & 什么都不用带，代码已经进 exe & 必须把 \code{.dll} 一起发布 \\
体积 & exe 变大，多个程序各自持有一份 & 多个程序共享一份 \code{.dll} \\
更新 & 改库就要重新编译发布主程序 & 替换 \code{.dll} 即可（ABI 不变的前提下） \\
排查难度 & 符号直接进 exe，问题相对好定位 & 多一层加载过程，缺文件、位数不符都会在启动时失败 \\
许可 & GPL 类库静态链接可能带来传染性义务 & 通常只要求随程序分发 \code{.dll} \\
\bottomrule
\end{longtable}
\subsection{用 dumpbin 查依赖}
\code{dumpbin} 不在普通命令行里，需要用 \ui{工具 > 命令行 > 开发者命令提示符} 打开，或从开始菜单启动 Developer Command Prompt。

\begin{lstlisting}[style=shellstyle]
:: 列出直接依赖的 DLL（不递归）
dumpbin /dependents MyApp.exe

:: 查看 PE 头，确认位数：x64 显示 8664 machine (x64)，x86 显示 14C machine (x86)
dumpbin /headers MyApp.exe | findstr machine

:: 查看导入表，能看到每个 DLL 被用到的具体函数
dumpbin /imports MyApp.exe

:: 检查某个第三方 DLL 的位数
dumpbin /headers helper.dll | findstr machine
\end{lstlisting}

\begin{tipbox}[提示]
\code{/dependents} 只列\textbf{直接}依赖，不会展开到第二层。要递归看整棵依赖树，用开源的 Dependencies 工具（Dependency Walker 的现代替代品），或者干脆在每个目录上手工跑一遍 \code{dumpbin}。
\end{tipbox}
\subsection{定位缺失 DLL 的排查路径}
\begin{enumerate}
  \item 先记下报错对话框里的 DLL 文件名，这是最直接的线索。
  \item \code{dumpbin /dependents} 列出直接依赖，逐个确认它们在目标机器上的位置。
  \item 用 Dependencies 递归展开，缺失的那一级会被标出来；它的 API Set 解析比老工具准确得多。
  \item 用 Process Monitor 过滤 \code{Process Name is MyApp.exe} 与 \code{Result is NAME NOT FOUND}，可以精确看到加载器依次找过哪些目录。
  \item 在 Visual Studio 里用 \ui{调试 > 窗口 > 模块} 查看实际加载到的模块全路径；\ui{工具 > 选项 > 调试 > 输出窗口} 里打开模块加载消息，也能在输出窗口看到加载过程。
  \item 注意最常见的陷阱：应用目录里同时存在 32 位与 64 位的同名 DLL，加载器挑中了错的那个，于是报 \code{0xc000007b}。
\end{enumerate}
\section{收集运行时依赖}
\subsection{Visual C++ 可再发行程序包}
可再发行程序包（Redistributable）就是把 CRT 装进系统目录的官方安装器，文件名按位数区分：
\begin{itemize}
  \item \code{vc_redist.x64.exe}——64 位程序用。
  \item \code{vc_redist.x86.exe}——32 位程序用（即使跑在 64 位 Windows 上也是它）。
  \item \code{vc_redist.arm64.exe}——ARM64 程序用。
\end{itemize}

要点：
\begin{itemize}
  \item \textbf{位数必须与你的程序一致}。64 位程序装 x86 版运行库没有任何作用，报错依旧。
  \item 安装包本体位于 VS 安装目录下的 \codeP{VC\textbackslash{}Redist\textbackslash{}MSVC\textbackslash{}<工具集版本>}，同级的 \codeP{x64\textbackslash{}Microsoft.VC<编号>.CRT} 目录里放的是供 app-local 部署的 CRT DLL 文件。
  \item 自 Visual Studio 2015 起，各代运行库保持二进制兼容并可以共存，装一个新的通常就能满足老程序；但老程序若需要 \code{VCRUNTIME140_1.dll}，就必须装较新的运行库。
  \item 通用 CRT（UCRT）从 Windows 10 起由系统提供；更早的系统需要额外安装系统更新（例如 KB2999226 一类），这也是“在 Win10 上好好的，到老系统就缺 \code{ucrtbase.dll}”的原因。
\end{itemize}
\subsection{随安装包分发与静默安装}
\begin{lstlisting}[style=shellstyle]
:: 静默安装运行库：/install 安装、/quiet 无界面、/norestart 不自动重启
vc_redist.x64.exe /install /quiet /norestart

:: 取返回值判断结果
echo %ERRORLEVEL%
\end{lstlisting}

常见返回值：

\begin{longtable}{L{2.2cm}L{4.2cm}L{8.2cm}}
\toprule
\textbf{返回值} & \textbf{含义} & \textbf{处理方式} \\
\midrule
0 & 安装成功 & 继续安装主程序 \\
1638 & 已存在其他版本，未安装本次版本 & 视为成功，继续 \\
3010 & 安装成功但需要重启 & 提示用户重启后生效 \\
1603 & 安装失败 & 检查权限、是否有更高版本阻止降级、日志中的具体错误 \\
\bottomrule
\end{longtable}

其他可用参数：\code{/passive} 显示进度条但不需要交互、\code{/repair} 修复、\code{/uninstall} 卸载、\code{/log <文件>} 输出日志。在安装包项目里，把这个 exe 作为普通文件一起打进去，再在安装流程末尾用批处理或自定义操作静默执行，是最省事的做法。
\subsection{app-local 部署的边界}
“把 CRT 的 DLL 直接拷到 exe 同目录”叫 app-local 部署，它能免掉安装步骤，但有几个前提：
\begin{itemize}
  \item 只有官方《可再发行代码》清单里列出的文件才允许随程序分发，不同 Visual Studio 版本与不同版本（Community、Professional、Enterprise）的授权范围并不相同，务必以随 VS 安装的 \code{Redist.txt} 和 Microsoft Learn 上的“可再发行代码”页面为准。
  \item app-local 的 CRT 不会随 Windows Update 被修补，安全更新要自己重新发布一次。
  \item 系统目录里若已有另一份 CRT 被其他进程加载，问题会更难排查。
\end{itemize}

结论：\textbf{对外正式发布优先用 \code{vc_redist} 安装器；绿色小工具或内网工具可以用 app-local。}
\section{生成安装包}
\subsection{三种打包方式对比}
\begin{longtable}{L{3.4cm}L{2.6cm}L{4.1cm}L{4.1cm}}
\toprule
\textbf{方案} & \textbf{产物} & \textbf{优点} & \textbf{局限} \\
\midrule
Visual Studio Installer Projects 扩展 & \code{.msi} 与 \code{setup.exe} & 直接集成在 VS 里，图形化添加快捷方式、注册表、文件类型关联 & 需要单独安装扩展；自定义能力有限；多语言与复杂升级场景支持一般 \\
MSIX 打包项目 & \code{.msix} & 安装卸载干净、支持增量更新与回滚、可上 Microsoft Store & 安装到受保护目录；文件与注册表被虚拟化；不能装服务与驱动；必须代码签名 \\
WiX Toolset & \code{.msi}（由 XML 描述） & 完全可脚本化，适合放进 CI，能表达复杂安装逻辑 & 学习曲线陡，要手写 XML \\
\bottomrule
\end{longtable}
\subsection{Installer Projects 最小制作流程}
\begin{enumerate}
  \item 打开 \ui{扩展 > 管理扩展}，搜索并安装 \code{Microsoft Visual Studio Installer Projects}，安装完成后重启 VS。
  \item \ui{文件 > 新建 > 项目}，搜索 "Setup Project"，选择安装项目模板。
  \item 视图切换到 \ui{文件系统 (File System)}，右键 \ui{应用程序文件夹} 节点，选择 \ui{添加 > 项目输出}，选中主项目；第三方 DLL 用 \ui{添加 > 文件} 补进去。
  \item 选中安装项目，在属性窗口里把 \textbf{TargetPlatform} 设为 \code{x64}，必须与 exe 的位数一致，否则会装进 \code{Program Files (x86)}。
  \item 右键 \ui{应用程序文件夹} 里的主输出，选择 \ui{创建主输出的快捷方式}，把生成的快捷方式拖到 \ui{用户桌面} 和 \ui{用户的“程序”菜单}。
  \item 在安装项目的属性里设置 \textbf{Manufacturer}、\textbf{Version}、\textbf{ProductName}；升级安装时要保持 \textbf{UpgradeCode} 不变，提升 \textbf{Version} 并把 \textbf{RemovePreviousVersions} 设为 \code{True}，否则会出现“已安装另一个版本”而无法覆盖。
  \item 执行 \ui{生成 > 生成解决方案}，在安装项目的 Release 输出目录里得到 \code{.msi} 与 \code{setup.exe}。
  \item 在\textbf{一台干净的虚拟机}上完整走一遍安装、运行、卸载。这是唯一真正有效的验证方式。
\end{enumerate}
\subsection{WiX Toolset 简介}
WiX 用 XML 描述安装内容：产品信息、目录树、组件与 GUID、快捷方式、注册表项、服务、自定义操作，然后由命令行工具编译成 MSI。它不做图形化设计器，换来的是可版本管理、可在 CI 里一条命令出包。

适合用 WiX 的信号：安装逻辑复杂（多组件可选安装、按条件写注册表）、需要无人值守构建、需要被内部部署系统调用。如果只是“把 exe 和几个 DLL 装进 Program Files 并建快捷方式”，Installer Projects 更快。
\section{其他发布形式}
\subsection{绿色免安装版}
做法很简单：用 \code{/MT} 加静态第三方库，或者把所需 DLL 与 exe 放进同一层目录，然后压缩成 zip。之所以“放同目录”最稳，是因为 Windows 加载器的搜索顺序是固定的（默认开启安全搜索模式）：

\begin{lstlisting}[style=textstyle]
DLL 搜索顺序（安全搜索模式，默认开启）
  1. 加载它的那个 exe/dll 所在目录  <- 唯一完全由你掌控的位置
  2. System32（32 位进程在 64 位系统上是 SysWOW64）
  3. Windows 目录
  4. 当前工作目录  <- 双击启动和命令行启动可能不同，别依赖
  5. PATH 环境变量中的各目录
\end{lstlisting}

进一步收紧可以用 \code{SetDefaultDllDirectories} 配合 \code{AddDllDirectory} / \code{LoadLibraryEx} 的搜索标志，只信任自己的目录，避免被同名 DLL 劫持。
\subsection{单文件与自包含}
“单文件发布”是 .NET 的概念（\code{PublishSingleFile}、\code{SelfContained}），C++ 没有对应的官方开关。C++ 里等价的做法是：静态 CRT + 静态第三方库 + 不依赖外部数据文件，产物天然就是一个 exe。

需要清楚的两点：\code{/MT} 只能消掉 CRT 依赖，消不掉对系统 DLL 的依赖；图标、字体、数据库之类的资源仍需外置，或者塞进资源段（在 \code{.rc} 里用 \code{RCDATA} 声明，运行时用 \code{FindResource} / \code{LoadResource} 读出来）。
\subsection{分发渠道一句话说明}
\begin{itemize}
  \item \textbf{Steam}：把 exe 与 \code{steam_api64.dll} 一起提交到 depot，并可用 Steamworks 的安装脚本（\code{installscript.vdf}）声明在首次运行前执行 VC++ 运行库安装程序。
  \item \textbf{内网服务器}：直接放 zip 加校验和，用 \code{certutil -hashfile MyApp.exe SHA256} 生成 SHA-256 供同事核对；也可以接内部的包管理源。
  \item \textbf{企业批量部署}：优先出 MSI，交给域策略或软件分发系统；绿色版作为兜底。
\end{itemize}
\section{版本信息与资源}
\subsection{添加 \code{.rc} 资源文件}
\begin{enumerate}
  \item 打开 \ui{项目 > 添加 > 资源}，选择 \textbf{Version}，新建后 VS 会生成 \code{<项目名>.rc} 与 \code{resource.h}。
  \item 再次打开 \ui{项目 > 添加 > 资源}，选择 \textbf{Icon}，导入 \code{.ico}（建议一个文件里包含 16、32、48、256 多档尺寸）。
  \item \code{.rc} 由资源编译器编译成 \code{.res}，再由链接器并入 exe，不需要你手动调用任何工具。
\end{enumerate}
\subsection{版本号字段}
以下是一个最小可用的版本资源写法：

\begin{lstlisting}[style=cppstyle]
#include "resource.h"

VS_VERSION_INFO VERSIONINFO
 FILEVERSION    1,2,3,0
 PRODUCTVERSION 1,2,3,0
 FILEFLAGSMASK  0x3fL
 FILEFLAGS      0x0L
 FILEOS         0x40004L
 FILETYPE       0x1L
 FILESUBTYPE    0x0L
BEGIN
    BLOCK "StringFileInfo"
    BEGIN
        BLOCK "080404b0"
        BEGIN
            VALUE "CompanyName",      "示例公司\0"
            VALUE "FileDescription",  "示例工具\0"
            VALUE "FileVersion",      "1.2.3.0\0"
            VALUE "InternalName",     "MyApp\0"
            VALUE "LegalCopyright",   "Copyright (C) 2026\0"
            VALUE "OriginalFilename", "MyApp.exe\0"
            VALUE "ProductName",      "MyApp\0"
            VALUE "ProductVersion",   "1.2.3.0\0"
        END
    END
    BLOCK "VarFileInfo"
    BEGIN
        VALUE "Translation", 0x804, 1200
    END
END
\end{lstlisting}

\begin{longtable}{L{6.0cm}L{8.8cm}}
\toprule
\textbf{字段} & \textbf{显示位置与作用} \\
\midrule
\code{FILEVERSION} / \code{PRODUCTVERSION} & 二进制形式，四个 16 位整数，用逗号分隔；供程序读取比较 \\
\code{FileVersion} / \code{ProductVersion} & 字符串形式，资源管理器“详细信息”里看到的就是它 \\
\code{CompanyName} & 公司名，安装包与“详细信息”页都会读 \\
\code{FileDescription} & 文件说明，任务管理器的进程描述列取这里 \\
\code{LegalCopyright} & 版权声明 \\
\code{OriginalFilename} & 原始文件名，改名后仍能追溯 \\
\code{080404b0} & 语言块：\code{0804} 为简体中文，\code{04b0} 为 Unicode 代码页 1200 \\
\bottomrule
\end{longtable}

\begin{warnbox}[注意]
\code{StringFileInfo} 的语言代码块必须与 \code{VarFileInfo} 里的 \code{Translation} 一致（上面是 \code{0x804, 1200}），否则部分工具读不到字符串。\code{FILEVERSION} 与 \code{FileVersion} 也要手工保持一致——它们不会互相推导。
\end{warnbox}
\subsection{图标与属性页设置}
\begin{itemize}
  \item exe 在资源管理器里显示的图标，取资源中 \textbf{ID 最小} 的那个图标，所以通常用 \code{IDI_ICON1}。
  \item 选中 \code{.rc} 文件后打开属性页，根节点就是 \textbf{资源}：\ui{资源 > 常规 > 区域性} 可以设置资源语言（中文简体为 \code{0x0804}）。不改的话默认是英文，多语言资源会显得混乱。
  \item 想让版本号跟随构建号自动递增，可以在 \ui{项目 > 属性 > 生成事件 > 生成前事件} 里调用脚本改写 \code{.rc}；也可以在 \ui{资源 > 常规 > 预处理器定义} 里注入宏，让 \code{.rc} 用 \code{#ifdef} 选择版本字符串。
\end{itemize}
\section{发布检查清单}
\textbf{配置与平台}
\begin{itemize}
  \item[$\square$] \ui{生成 > 配置管理器} 里所有项目都是 Release，平台都是 x64（或与目标一致）。
  \item[$\square$] \ui{项目 > 属性 > C/C++ > 命令行} 中不再出现 \code{/Od}、\code{/RTC1}、\code{/MDd}、\code{/MTd}。
  \item[$\square$] 全解决方案执行过 \ui{重新生成解决方案}，没有沿用 Debug 的中间文件。
  \item[$\square$] \ui{项目 > 属性 > 常规 > 平台工具集} 与 \ui{Windows SDK 版本} 在开发机与构建机上一致。
\end{itemize}

\textbf{依赖}
\begin{itemize}
  \item[$\square$] \code{dumpbin /dependents MyApp.exe} 的输出里\textbf{没有} \code{MSVCP140D.dll}、\code{VCRUNTIME140D.dll}、\code{ucrtbased.dll} 之类的调试版运行库。
  \item[$\square$] 所有第三方 DLL 都是 Release 版、位数正确、CRT 设置与本项目一致。
  \item[$\square$] 明确记录：目标机器需要装哪个可再发行程序包（x86 还是 x64）。
  \item[$\square$] 若有数据文件、字体、配置模板，也一并纳入安装包或目录结构。
\end{itemize}

\textbf{测试}
\begin{itemize}
  \item[$\square$] 在干净虚拟机上完成安装、首次运行、卸载全流程。
  \item[$\square$] 在\textbf{没有装过开发工具}的机器上验证一遍，重点看启动阶段是否缺 DLL。
  \item[$\square$] 中文路径、含空格用户名、非管理员账户下各跑一次。
  \item[$\square$] 若是数值计算类程序，比较 Release 与 Debug 的关键输出是否一致（防止 \code{/fp:fast} 引入偏差）。
\end{itemize}

\textbf{版本与资源}
\begin{itemize}
  \item[$\square$] exe 的“详细信息”页里能看到正确的版本号、公司名、版权。
  \item[$\square$] 图标在资源管理器、任务栏、任务管理器中显示正常。
  \item[$\square$] 版本号与安装包的 \textbf{Version} 属性一致。
\end{itemize}

\textbf{签名与体积}
\begin{itemize}
  \item[$\square$] 用 \code{signtool sign /fd SHA256 /tr <时间戳服务器> /td SHA256 MyApp.exe} 完成签名，并用 \code{signtool verify /pa /v MyApp.exe} 验证。
  \item[$\square$] 签名带时间戳，这样证书到期后旧版本的签名依然有效。
  \item[$\square$] 体积已用 \code{/Gy} + \code{/OPT:REF} + \code{/OPT:ICF} 优化，且链接进来的库确实都被用到。
  \item[$\square$] 压缩包与安装包文件名里带版本号，便于回滚。
\end{itemize}
\section{发布后常见故障}
\subsection{缺少 \code{VCRUNTIME140.dll} 一类}
\textbf{表现}：弹窗“由于找不到 VCRUNTIME140.dll，无法继续执行代码。重新安装程序可能会解决此问题。”，常见变体还有 \code{MSVCP140.dll}、\code{VCRUNTIME140_1.dll}、\code{ucrtbase.dll}。

\textbf{原因}：程序用 \code{/MD} 编译，但目标机器没装对应位数的可再发行程序包；或者装的是 x86 版而程序是 x64；或者运行库太旧，缺 \code{VCRUNTIME140_1.dll}。

\textbf{处理}：
\begin{enumerate}
  \item 立即解决：在目标机运行 \code{vc_redist.x64.exe /install /quiet /norestart}。
  \item 根治：把运行库作为安装包的先决条件自动安装。
  \item 另一种根治：改用 \code{/MT} 重新发布，代价是体积变大、安全更新要自己跟。
\end{enumerate}
\subsection{应用程序无法正常启动 \code{0xc000007b}}
\code{0xC000007B} 是 \code{STATUS_INVALID_IMAGE_FORMAT}，意思是“这个映像文件的格式不对”。最常见的原因是\textbf{位数不匹配}——64 位进程加载了 32 位 DLL，或反过来；其次是 DLL 文件被截断或损坏、同名文件被不同位数的版本互相覆盖。

\begin{lstlisting}[style=textstyle]
x64 的 MyApp.exe --> 加载 helper.dll
                     |
                     +-- helper.dll 是 32 位    --> 0xC000007B（映像格式无效）
                     +-- helper.dll 不存在      --> 0xC0000135（找不到 DLL）
                     +-- helper.dll 初始化失败  --> 0xC0000142（DLL 初始化失败）
\end{lstlisting}

排查方法：

\begin{lstlisting}[style=shellstyle]
:: 先确认主程序位数，再逐个确认可疑 DLL 的位数
dumpbin /headers MyApp.exe  | findstr machine
dumpbin /headers helper.dll | findstr machine
\end{lstlisting}

两者的 \code{machine} 值必须一致：x64 显示 \code{8664 machine (x64)}，x86 显示 \code{14C machine (x86)}。顺带记住另两个高频代码：\code{0xC0000135} 是找不到 DLL，\code{0xC0000142} 是 DLL 初始化失败（常见于被杀软拦截或它的依赖没满足）。
\subsection{DLL 地狱的标准排查步骤}
\begin{enumerate}
  \item 抄下完整报错文案，记录故障机器的 Windows 版本与位数。
  \item \code{dumpbin /dependents MyApp.exe} 列出直接依赖，逐个确认在目标机器上的存在性与位数。
  \item 用 Dependencies 递归展开依赖树，找出断链的那一级。
  \item 用 Process Monitor 过滤进程名与 \code{NAME NOT FOUND}，看加载器实际搜索过哪些路径。
  \item 检查 PATH、当前工作目录，以及 System32 里是否存在同名 DLL。
  \item 在“能跑起来”的机器上用 \ui{调试 > 窗口 > 模块} 导出已加载模块清单，与故障机逐条对比。
  \item 确认三件事：位数一致、CRT 配置一致（同为 \code{/MD} 或同为 \code{/MT}）、没有 Debug 版 CRT 混入。
  \item 修完后在干净虚拟机上复测一遍，不要只在开发机上确认。
\end{enumerate}
\section{小结}
\begin{itemize}
  \item Debug 与 Release 的差异是一整组开关的差异：优化、宏、运行库、调试信息、增量链接，切换时必须全解决方案一致。
  \item \code{/O2}、\code{/GL} + \code{/LTCG}、\code{/arch:AVX2} 能显著提速，但分别带来调试困难、构建变慢、以及在不支持该指令集的 CPU 上崩溃的风险。
  \item \code{/MD} 依赖可再发行程序包但体积小、跨模块安全；\code{/MT} 免依赖但体积大且跨模块传 STL 危险；同一程序内混用会导致堆损坏，DLL 层面的不一致链接期发现不了。
  \item 静态库无部署成本但更新要重编；动态库必须随程序带上 \code{.dll}，用 \code{dumpbin /dependents} 与 Dependencies 工具确认依赖完整。
  \item 运行库按位数分发，\code{vc_redist.x64.exe /install /quiet /norestart} 可静默安装，返回值 0、1638、3010、1603 要分别处理。
  \item 安装包三条路线：Installer Projects 最快、MSIX 最干净但限制多、WiX 最适合 CI。
  \item 版本信息写进 \code{.rc}，exe 图标取资源中 ID 最小的图标，属性页里的 \ui{资源 > 常规 > 区域性} 决定资源语言。
  \item 发布后两类高频故障是缺运行库和 \code{0xc000007b}（位数不匹配），排查顺序都是先确认位数、再确认依赖树。
\end{itemize}
\section{练习题}
\begin{enumerate}
  \item 把一个控制台项目从 \code{/MD} 改成 \code{/MT}，用 \code{dumpbin /dependents} 对比改动前后的依赖列表，并解释每一行的来源。
  \item 用 \code{/arch:AVX2} 编译一个程序，分别在支持与不支持 AVX2 的机器（或限制 CPU 特性的虚拟机）上运行，记录现象与事件查看器中的异常代码。
  \item 给一个项目添加版本资源，使 exe 的“详细信息”页显示公司名与版本号 \code{1.0.0.0}，然后通过代码读取并打印 \code{FILEVERSION} 与 \code{FileVersion}，说明二者的区别。
  \item 制作一个最小 MSI：包含主程序 exe 与一个第三方 DLL，安装后在开始菜单创建快捷方式，并在干净虚拟机上完成安装与卸载测试。
  \item 故意把 32 位的 DLL 放进 64 位程序的目录，复现 \code{0xc000007b}，写出从报错到定位的完整排查过程。
  \item 写一个批处理，在安装主程序前静默安装 \code{vc_redist.x64.exe}，并正确处理返回值 0、1638、3010、1603。
\end{enumerate}
\section{相关笔记}
\begin{itemize}
  \item \textbf{本教程}：\textbf{Visual Studio 2026 使用教程}
  \item \textbf{上一章}：\textbf{单元测试与代码质量}
  \item \textbf{下一章}：\textbf{常见问题与效率技巧}
  \item \textbf{CMake}：\textbf{CMake 使用指南}
  \item \textbf{Git 命令}：\textbf{常用 Git 命令}
\end{itemize}
